\textbf{11.} Пусть $C$ — случайный язык, который для каждого $n$ с вероятностью $\frac12$ не содержит слов длины $n$, а с вероятностью $\frac12$ содержит ровно одно слово длины $n$ (равновероятно по всем словам такой длины). Докажите, что с высокой вероятностью $\p^C \neq \np^C$

\textbf{Решение:} Так как $\p \subset \np$, то для любого оракула $C$ $\p^C \subset \np^C$. Для того, чтобы показать неравенство, нужно предоставить язык $U$, лежащий в $\np^C$, но не в $\p^C$.

Зафиксируем $C$. Рассмотрим $U = \{1^n \mid$ в $C$ существует слово длины $n\}$ (можно сказать, что мы упростили язык $C$ до унарного). 

\textbf{Утверждение 1:} $U \in \np^C$.  Построим недетерминированную машину Тьюринга $M(x)$, которая может обращаться к языку $C$, для доказательства этого утверждения. Пусть машине $M$ пришло слово $1^n$, тогда недетерминизм позволяет нам перебрать все слова длины $n$ за полиномиальное время (если вам непонятно, как, то обратитесь к определению $\np$), а обращение к оракулу позволит за $O(1)$ проверить их принадлежность к $C$. Если монетка выпала решкой, и какое-то случайное слово длины $n$ есть в языке $C$, то мы его переберем, проверим; скажем, что $1^n \in U$, и будем правы. Иначе, если монетка выпала орлом и такого слова нет в $C$, то ответим $1^n \not\in U$. Это и доказывает принадлежность языка к $\np^C$.

\textbf{Утверждение 2:} $U \not \in \p^C$.  Будут приведены лишь косвенные рассуждения, объясняющие, почему это так. Предположим, что это не так и найдется детерминированная машина Тьюринга $M$, работающая за полиномиальное время и умеющая обращаться к языку $C$. Тогда время её работы ограничено неким полиномом $Dn^k,~ D,k>0$. Всего вариантов слов длины $n$ - $2^n$, и каждое слово равновероятно. Хочется сказать, что только по длине слова машина $M$ не может проверить присутствие или отсутствие слова такой длины в $C$ эффективнее, чем перебором всех слов. Но ввиду того, что слов эспоненциально много, а машина работает полиномиально долго, то перебор покроет максимум $\frac{Dn^k}{2^n}$ слов. Эта доля со стремлением $n$ к бесконечности стремится к 0, поэтому вероятность, что $U \in \p^C$ стремится к 0, а значит хотя бы с вероятностью 1 утверждение 2 верно и $\p^C \neq \np^C$. 
